\section{Tempi, efficienza e verifica del parallelismo}\label{app:parallel-results}

\subsection{Definizioni e lettura dei rapporti}
Le metriche $A_a(p)$, $S_a(p)$ ed $E_a(p)$ sono definite dall'Eq.~\eqref{eq:speedup-efficiency}. Nel caso dei backend sincroni le misure si riferiscono a una stessa traiettoria, verificata mediante i campi registrati per seme e lavoro. Le architetture asincrona e a isole modificano invece il processo stocastico: il confronto temporale rispetto alla seriale ha significato di \emph{runtime ratio} fra run complete. La selezione del framework usa il wall-clock, mentre il timer nativo distingue il costo del solver da quello dell'invocazione.

\subsection{Curve complementari dei backend sincroni}
Le Figure~\ref{fig:new-sync-serial-rastrigin} e~\ref{fig:new-sync-serial-diffusion} riportano i rapporti rispetto alla seriale nel campione di selezione pilot; le Figure~\ref{fig:new-sync-ext-rast} e~\ref{fig:new-sync-ext-diff} completano le curve di strong scaling ed efficienza presenti nel testo. Obiettivi e dimensioni non sono aggregati perché hanno granularità di lavoro differente.

\begin{figure}[H]\centering
\includegraphics[width=0.93\textwidth,height=0.72\textheight,keepaspectratio]{02_parallel/selection/fig_serial_speedup_rastrigin_d30.pdf}
\caption{Risultati dei backend sincroni nel pilot: accelerazione rispetto alla seriale su Rastrigin a $D=30$. Il riferimento temporale e la semantica delle traiettorie sono quelli definiti dall'Eq.~\eqref{eq:speedup-efficiency}.}\label{fig:new-sync-serial-rastrigin}
\end{figure}
\begin{figure}[H]\centering
\includegraphics[width=0.93\textwidth,height=0.72\textheight,keepaspectratio]{02_parallel/selection/fig_serial_speedup_rastrigin_d100.pdf}
\caption{Risultati dei backend sincroni nel pilot: accelerazione rispetto alla seriale su Rastrigin a $D=100$. Il riferimento temporale e la semantica delle traiettorie sono quelli definiti dall'Eq.~\eqref{eq:speedup-efficiency}.}\label{fig:new-sync-serial-rastrigin-extra}
\end{figure}
\begin{figure}[H]\centering
\includegraphics[width=0.93\textwidth,height=0.72\textheight,keepaspectratio]{02_parallel/selection/fig_serial_speedup_inverse_diffusion_d30.pdf}
\caption{Risultati dei backend sincroni nel pilot: accelerazione rispetto alla seriale su diffusione inversa a $D=30$. Il riferimento temporale e la semantica delle traiettorie sono quelli definiti dall'Eq.~\eqref{eq:speedup-efficiency}.}\label{fig:new-sync-serial-diffusion}
\end{figure}
\begin{figure}[H]\centering
\includegraphics[width=0.93\textwidth,height=0.72\textheight,keepaspectratio]{02_parallel/selection/fig_serial_speedup_inverse_diffusion_d100.pdf}
\caption{Risultati dei backend sincroni nel pilot: accelerazione rispetto alla seriale su diffusione inversa a $D=100$. Il riferimento temporale e la semantica delle traiettorie sono quelli definiti dall'Eq.~\eqref{eq:speedup-efficiency}.}\label{fig:new-sync-serial-diffusion-extra}
\end{figure}
\begin{figure}[H]\centering
\includegraphics[width=0.93\textwidth,height=0.72\textheight,keepaspectratio]{02_parallel/selection/fig_strong_speedup_rastrigin_d100.pdf}
\caption{Risultati dei backend sincroni nel pilot: strong scaling su Rastrigin a $D=100$. Il riferimento temporale e la semantica delle traiettorie sono quelli definiti dall'Eq.~\eqref{eq:speedup-efficiency}.}\label{fig:new-sync-ext-rast}
\end{figure}
\begin{figure}[H]\centering
\includegraphics[width=0.93\textwidth,height=0.72\textheight,keepaspectratio]{02_parallel/selection/fig_efficiency_rastrigin_d100.pdf}
\caption{Risultati dei backend sincroni nel pilot: efficienza interna su Rastrigin a $D=100$. Il riferimento temporale e la semantica delle traiettorie sono quelli definiti dall'Eq.~\eqref{eq:speedup-efficiency}.}\label{fig:new-sync-ext-rast-extra}
\end{figure}
\begin{figure}[H]\centering
\includegraphics[width=0.93\textwidth,height=0.72\textheight,keepaspectratio]{02_parallel/selection/fig_strong_speedup_inverse_diffusion_d30.pdf}
\caption{Risultati dei backend sincroni nel pilot: strong scaling su diffusione inversa a $D=30$. Il riferimento temporale e la semantica delle traiettorie sono quelli definiti dall'Eq.~\eqref{eq:speedup-efficiency}.}\label{fig:new-sync-ext-diff}
\end{figure}
\begin{figure}[H]\centering
\includegraphics[width=0.93\textwidth,height=0.72\textheight,keepaspectratio]{02_parallel/selection/fig_efficiency_inverse_diffusion_d30.pdf}
\caption{Risultati dei backend sincroni nel pilot: efficienza interna su diffusione inversa a $D=30$. Il riferimento temporale e la semantica delle traiettorie sono quelli definiti dall'Eq.~\eqref{eq:speedup-efficiency}.}\label{fig:new-sync-ext-diff-extra}
\end{figure}

\subsection{Studio precedente del costo sintetico}
La campagna a carico controllato utilizza Sphere con lavoro deterministico aggiuntivo: l'obiettivo restituisce Sphere più $10^{-30}$ volte il valore assoluto di un accumulatore aggiornato mediante \texttt{fma}. Il numero di operazioni controlla il costo per FE e viene calibrato separatamente in ogni sessione. La configurazione impiega $N=64$, $D=30$, vicinato Global, casualità scalare e $h=1$. Ai costi nominali $0$, $10$, $100$, $1000\,\mu$s corrispondono rispettivamente $K=31$, $31$, $14$, $1$ aggiornamenti, cioè $2\,048$, $2\,048$, $960$, $128$ FE per run. Ogni lotto contiene otto semi e due ripetizioni per sessione; tre sessioni producono 576 lotti e 4\,608 esecuzioni individuali.

Il controllo delle esecuzioni sincrone confronta esattamente i campi osservati \texttt{best\_value}, \texttt{best\_position}, \texttt{best\_history} e le ascisse FE. Nei lotti con repliche concorrenti, la somma dei timer nativi misura il lavoro; il confronto di throughput usa invece il cronometro esterno, inclusi creazione e terminazione dei processi. Questo esperimento fornisce una misura complementare della sensibilità al costo, separata dal pilot che sceglie il framework in base a qualità e tempo.

\subsection{Rapporti temporali e modello di overhead}
\small
\begin{longtable}{@{}llrrrrrr@{}}
\caption{Tempi nativi nella sessione 3. MS: master--slave; PP: particle-parallel. Costo zero indica assenza del lavoro sintetico aggiuntivo. I valori sono mediane di due rapporti appaiati; $H$ è definito a un worker. Piccoli valori negativi riflettono la variabilità temporale.}\label{tab:parallel-full}\\
\toprule Dinamica & Costo $\mu$s & Backend & $p$ & $A(p)$ & $S(p)$ & $E$ \% & $H$ \%\\\midrule\endfirsthead
\toprule Dinamica & Costo $\mu$s & Backend & $p$ & $A(p)$ & $S(p)$ & $E$ \% & $H$ \%\\\midrule\endhead
\bottomrule\endfoot
A & 0 & MS & 1 & 0.835 & 1.000 & 100.0 & 19.74\\
A & 0 & MS & 2 & 0.818 & 0.980 & 49.0 & ---\\
A & 0 & MS & 4 & 0.787 & 0.943 & 23.6 & ---\\
A & 0 & MS & 8 & 0.669 & 0.800 & 10.0 & ---\\
A & 0 & PP & 1 & 0.798 & 1.000 & 100.0 & 25.25\\
A & 0 & PP & 2 & 1.132 & 1.418 & 70.9 & ---\\
A & 0 & PP & 4 & 1.343 & 1.682 & 42.0 & ---\\
A & 0 & PP & 8 & 0.911 & 1.140 & 14.3 & ---\\
A & 10 & MS & 1 & 0.962 & 1.000 & 100.0 & 3.99\\
A & 10 & MS & 2 & 1.676 & 1.743 & 87.2 & ---\\
A & 10 & MS & 4 & 2.552 & 2.654 & 66.3 & ---\\
A & 10 & MS & 8 & 2.938 & 3.055 & 38.2 & ---\\
A & 10 & PP & 1 & 0.937 & 1.000 & 100.0 & 6.68\\
A & 10 & PP & 2 & 1.758 & 1.876 & 93.8 & ---\\
A & 10 & PP & 4 & 2.885 & 3.079 & 77.0 & ---\\
A & 10 & PP & 8 & 3.598 & 3.838 & 48.0 & ---\\
A & 100 & MS & 1 & 1.032 & 1.000 & 100.0 & -3.00\\
A & 100 & MS & 2 & 1.992 & 1.929 & 96.5 & ---\\
A & 100 & MS & 4 & 3.728 & 3.610 & 90.3 & ---\\
A & 100 & MS & 8 & 5.896 & 5.710 & 71.4 & ---\\
A & 100 & PP & 1 & 1.029 & 1.000 & 100.0 & -2.65\\
A & 100 & PP & 2 & 2.002 & 1.946 & 97.3 & ---\\
A & 100 & PP & 4 & 3.671 & 3.564 & 89.1 & ---\\
A & 100 & PP & 8 & 5.357 & 5.230 & 65.4 & ---\\
A & 1000 & MS & 1 & 1.001 & 1.000 & 100.0 & -0.12\\
A & 1000 & MS & 2 & 1.937 & 1.935 & 96.7 & ---\\
A & 1000 & MS & 4 & 3.636 & 3.632 & 90.8 & ---\\
A & 1000 & MS & 8 & 5.637 & 5.630 & 70.4 & ---\\
A & 1000 & PP & 1 & 0.999 & 1.000 & 100.0 & 0.10\\
A & 1000 & PP & 2 & 1.935 & 1.936 & 96.8 & ---\\
A & 1000 & PP & 4 & 3.690 & 3.694 & 92.3 & ---\\
A & 1000 & PP & 8 & 6.273 & 6.279 & 78.5 & ---\\
D & 0 & MS & 1 & 0.784 & 1.000 & 100.0 & 27.60\\
D & 0 & MS & 2 & 0.719 & 0.918 & 45.9 & ---\\
D & 0 & MS & 4 & 0.674 & 0.860 & 21.5 & ---\\
D & 0 & MS & 8 & 0.558 & 0.712 & 8.9 & ---\\
D & 0 & PP & 1 & 0.682 & 1.000 & 100.0 & 46.63\\
D & 0 & PP & 2 & 0.856 & 1.254 & 62.7 & ---\\
D & 0 & PP & 4 & 0.905 & 1.328 & 33.2 & ---\\
D & 0 & PP & 8 & 0.579 & 0.851 & 10.6 & ---\\
D & 10 & MS & 1 & 0.967 & 1.000 & 100.0 & 3.46\\
D & 10 & MS & 2 & 1.745 & 1.806 & 90.3 & ---\\
D & 10 & MS & 4 & 2.943 & 3.045 & 76.1 & ---\\
D & 10 & MS & 8 & 3.483 & 3.604 & 45.0 & ---\\
D & 10 & PP & 1 & 0.961 & 1.000 & 100.0 & 4.03\\
D & 10 & PP & 2 & 1.804 & 1.876 & 93.8 & ---\\
D & 10 & PP & 4 & 3.047 & 3.171 & 79.3 & ---\\
D & 10 & PP & 8 & 3.520 & 3.662 & 45.8 & ---\\
D & 100 & MS & 1 & 1.002 & 1.000 & 100.0 & -0.16\\
D & 100 & MS & 2 & 1.943 & 1.940 & 97.0 & ---\\
D & 100 & MS & 4 & 3.636 & 3.630 & 90.7 & ---\\
D & 100 & MS & 8 & 5.688 & 5.676 & 71.0 & ---\\
D & 100 & PP & 1 & 1.000 & 1.000 & 100.0 & -0.04\\
D & 100 & PP & 2 & 1.954 & 1.953 & 97.7 & ---\\
D & 100 & PP & 4 & 3.615 & 3.614 & 90.4 & ---\\
D & 100 & PP & 8 & 5.624 & 5.619 & 70.2 & ---\\
D & 1000 & MS & 1 & 0.999 & 1.000 & 100.0 & 0.14\\
D & 1000 & MS & 2 & 1.939 & 1.942 & 97.1 & ---\\
D & 1000 & MS & 4 & 3.652 & 3.657 & 91.4 & ---\\
D & 1000 & MS & 8 & 5.654 & 5.662 & 70.8 & ---\\
D & 1000 & PP & 1 & 1.000 & 1.000 & 100.0 & 0.03\\
D & 1000 & PP & 2 & 1.930 & 1.931 & 96.5 & ---\\
D & 1000 & PP & 4 & 3.640 & 3.641 & 91.0 & ---\\
D & 1000 & PP & 8 & 5.978 & 5.981 & 74.8 & ---\\
\end{longtable}\normalsize

Il modello dell'Eq.~\eqref{eq:overhead-model} è stimato con minimi quadrati non negativi per backend e dinamica, utilizzando le prime due sessioni. Il termine $aB$ descrive il costo seriale proporzionale al numero delle FE, $bB\widehat\tau_f/p$ il contributo delle valutazioni distribuite e $cp$ un overhead dipendente dal numero dei worker. La terza sessione è utilizzata esclusivamente per la verifica.
\begin{table}[H]\centering\small
\caption{Coefficienti non negativi del modello stimati sulle sessioni 1--2. La risposta è il tempo nativo totale delle otto run del lotto.}\label{tab:model}
\begin{tabular}{@{}llrrr@{}}
\toprule Dinamica & Backend & $a$ ($\mu$s/FE) & $b$ & $c$ (ms/worker)\\\midrule
A & MS & 0.453 & 0.9962 & 5.220\\
A & PP & 0.334 & 1.0020 & 4.899\\
D & MS & 0.000 & 0.9998 & 5.843\\
D & PP & 0.000 & 1.0006 & 6.020\\
\bottomrule\end{tabular}
\end{table}

\begin{table}[H]\centering\small
\caption{Verifica di tutte le 16 decisioni del modello nella sessione di hold-out. L'oracolo minimizza la mediana dei tempi osservati entro backend.}\label{tab:regret}
\begin{tabular}{@{}llrrrr@{}}
\toprule Dinamica & Costo $\mu$s & Backend & Scelta & Oracolo & Regret \%\\\midrule
A & 0 & MS & 1 & 1 & 0.0\\
A & 0 & PP & 1 & 4 & 68.2\\
A & 10 & MS & 4 & 8 & 14.9\\
A & 10 & PP & 8 & 8 & 0.0\\
A & 100 & MS & 8 & 8 & 0.0\\
A & 100 & PP & 8 & 8 & 0.0\\
A & 1000 & MS & 8 & 8 & 0.0\\
A & 1000 & PP & 8 & 8 & 0.0\\
D & 0 & MS & 1 & 1 & 0.0\\
D & 0 & PP & 1 & 4 & 32.5\\
D & 10 & MS & 4 & 8 & 18.4\\
D & 10 & PP & 4 & 8 & 15.5\\
D & 100 & MS & 8 & 8 & 0.0\\
D & 100 & PP & 8 & 8 & 0.0\\
D & 1000 & MS & 8 & 8 & 0.0\\
D & 1000 & PP & 8 & 8 & 0.0\\
\bottomrule\end{tabular}
\end{table}

La mediana dell'errore relativo assoluto di previsione è $7.15\%$ (analitico master--slave), $6.42\%$ (analitico particle-parallel), $6.92\%$ (discreto master--slave) e $6.17\%$ (discreto particle-parallel). La scelta coincide con il worker più rapido in 11 celle su 16. La misura di regret del modello documenta il costo di una scelta non ottimale, particolarmente visibile sui carichi più economici. In presenza di lavoro equivalente e costi dell'ordine del millisecondo, la CPSO analitica particle-parallel a otto worker raggiungeva un'accelerazione $A(8)=6.27$ e un'efficienza interna del $78.5\%$.
\FloatBarrier
